\section{Introducción}\label{sec:introduccion}
El estudio de las ecuaciones diferenciales parciales (EDPs) ha sido uno de los pilares fundamentales en el desarrollo de diversas áreas de las matemáticas aplicadas, desde la física teórica hasta la ingeniería. En particular, las EDPs no lineales juegan un papel crucial en la modelización de fenómenos complejos, como la propagación de ondas en medios no lineales o la turbulencia en fluidos.\\

Este trabajo se centra en la ecuación no linealidad modificada de Zakharov-Kusnetsov. Nuestro objetivo principal es estudiar el buen planteamiento de dicha ecuación en el contexto de los espacios de Sobolev periódicos abordandolo desde la perspectiva de las series de Fourier, algunas herramientas de análisis funcional y teoria de la medida.

A lo largo del documento, emplearemos la teoría de la transformada de Fourier y técnicas de análisis funcional para demostrar la existencia y unicidad de soluciones en estos espacios. Adicionalmente, exploraremos la convergencia de la serie de Fourier y algunas de las propiedades que se cumplen en los espacios de Sobolev, con el fin de formar un marco teórico sólido que nos permita solucionar nuestro problema.

El trabajo está organizado de la siguiente manera:
\begin{itemize}
    \item En la Sección 2 se introducen los preliminares necesarios, incluyendo los espacios vectoriales, la transformada de Fourier y algunas propiedades sobre el espacio de Sobolev.
    \item En la Sección 3 se presenta el resultado principal sobre el buen planteamiento de la ecuación no linealidad modificada de Zakharov-Kusnetsov.
    \item Y por último se encontrará la sección de referencias.
\end{itemize}
Finalmente cabe recalcar que este trabajo fue totalmente motivado, inspirado y fundamentado en el curso de Ecuaciones Diferenciales Parciales de la Universidad Nacional de Colombia, las notas de clase y los libros \cite{MR2597943} y \cite{MR3243734}.